
\chapter{Theoretical Framework}
\label{chap:two}

This thesis conducts analysis under the constructivist paradigm and uses \ac{OS} theory as the understanding for analysis of political speech acts by politicians in \ac{BiH}. Using this framework allows for a nuanced conceptualization of how political speech acts shape and are shaped by national identities, state actions, ethnic tensions, and other societal factors of state collectives.

\section{Social Constructivism}
\label{sec:constructivism}

Social constructivism provides an \ac{IR} framework that is largely based on the concept of intersubjectivity---shared meaning and understanding---in society and how various actors influence these meanings and understandings. The primary assumption that constitutes this intersubjectivity is ``that people act toward objects, including other actors, on the basis of the meanings that the objects have for them'' \citep[pg.~396-397]{wendt1992}. Furthermore, \citet[pg.~404.405]{wendt1992} highlights that actions taken before, during, and after an interaction between actors shapes the meaning of said interaction and how it is understood between them. Despite the relative simplicity of these concepts, they are essential to the constructivist paradigm and allow for the collectivization of meanings and understandings. Building upon these fundamental assumptions, actors and groups of actors are able to create complex societies, interactions, and institutions that constitute modern society, particularly that of governments, states, and the international global order \citep[pg.~401-402]{wendt1992}.

\subsection{Identity}
\label{subsec:identity}

Identity is fundamental to actors or groups of actors, but they are not limited to one particular identity, but rather multiple identities that are situation dependent. As \citet[pg.~397]{wendt1992} explains, ``Actors acquire identities---relatively stable, role-specific understandings and expectations about self---by participating in such collective meanings.'' The primary parts of identity, therefore, are an understanding of self and collective meaning under certain circumstances. This can be seen in \ac{BiH} when discussing \textit{Bosnian} identity and ones \textit{ethnic} identity. Actors often identify as \textit{Bosnian} when interacting with outsiders, but they identify as a particular \textit{ethnic} identity when interacting with other \textit{Bosnians} or members of their collective \textit{ethnic} identities. Contextuality determines many aspects of identity due to its intersubjectivity, and this allows identities to change as context may also change \citep[pg.~397-398]{wendt1992}.

Due to the reliance on roles, identities often overlap, such as someone being both a father and a manager. These two identities are not mutually exclusive, and may be beneficial or detrimental to each other under certain circumstances, but they are both reflections of self expectations and self understanding of one's roles \citep[pg.~397]{wendt1992}. When applying this to collective identities, such as national identities, ethnic identities, etc., there is a certain role that members of the collective assign to that identity, but also to other identities of a similar collective level \citep[pg.~280]{weldes1996}. These identities require other collective identities of a similar level---nations require other nations---because these identities are ``mutually constitutive'' \citep[as cited in \citealt{wendt1992}, pg.~399]{giddens1979}. Not only are identities formed by internal identification of \textit{self}, but also from external identification that is interpreted by the actor or group of actors who belong to that identity \citep[pg.~10-11]{todorova2009}.

On a national level, identities are formed by the collectivization of people that understand some belonging to a nation, but the identities themselves are shaped by the actions of statesmen \citep[pg.~281]{weldes1996}. National identity is still a collective understanding by the members of the nation, but the norms that form the institution of government provide statesmen with the agency to act and interpret meaning on behalf of the nation \citep[pg.~281]{weldes1996}. This agency provides the grounding for understanding the function that political speech acts in the formation, change, or dissolution of national identities.


\subsection{Interests}
\label{subsec:interest}

Following the formation of identities, actors form interests that correspond to one or more identities \citep[pg.~398]{wendt1992}. These interests are highly dependent on situations and are motivated by anticipated outcomes of each situation \citep[pg.~398]{wendt1992}, but these interests also serve to reproduce the understood meaning of an actor's identity \citep[pg.~347]{mitzenOntologicalSecurityWorld2006}. These interests come both from the actor's understood meaning of identity, but also as a societal understanding that has become institutionalized, as all social interactions and understanding are inherently intersubjective (\citealp[pg.~398-399]{wendt1992}; \citealp[pg.~280-281]{weldes1996}).

Similar to national identities, national interests are also primarily formed by statesmen in alignment with their understanding of the national identity, the situation, and potentially their personal interests \citep[pg.~280-282]{weldes1996}. \citet[pg.~280]{weldes1996} critiqued \citepos{wendt1992} interpretation of states as unitary actors, and posed that ``the meanings which objects, events and actions have for `states' are necessarily the meanings they have for those individuals who act in the name of the state.'' This interpretation does not mean that individual citizens of a state or nation do not impact national interests---citizens elect these statesmen in democracies, and protest provide a voice to contest actions made by statesmen---but it puts the primary formation of national interests in statesmen's hands.


\section{Ontological Security}
\label{sec:ontological_security}

The use of \ac{OS} in \ac{IR} research has increased significantly since it was introduced to the field by \citet{kinnvall2004} and \citet{mitzenOntologicalSecurityWorld2006}, who extended the original application from sociology \citep{giddens1991} to states and international actors. \ac{OS} relies on actors maintaining consistent identities, routines, trust, and other cognitive and social structures, which are often considered making up the \textit{security of being} \citep[pg.~240]{kinnvall2020a}. \citet{giddens1991} focused on the \ac{OS} of individuals, particularly the experiences of individuals and how they shape one's understanding of themselves. As \citet[pg.~352]{mitzenOntologicalSecurityWorld2006} poses, a similar approach to analyzing \textit{the security of being} can be applied to states as collective units formed by societal interactions and understanding. \ac{OS} in \ac{IR} expands upon the paradigm of social constructivism by applying \citepos{giddens1991} understanding of routines, trust, self to the state level \citep{mitzenOntologicalSecurityWorld2006, kinnvall2004}. This expansion constructed an analytical framework that provides an understanding of \ac{OS} and insecurity, as well as the ability of states to form interests and take actions with an inherent interest to maintain \ac{OS} and reduce uncertainty \citep{mitzenOntologicalSecurityWorld2006,kinnvall2004}.

\subsection{Identity}
\label{subsec:ont_identity}

Similar to traditional social constructivism, identities motivate actors' interests, actions, and understandings, but \ac{OS} expands the role of identity and poses that actors are motivated to maintain stable identities to feel ontologically secure \citep[pg.~344-345]{mitzenOntologicalSecurityWorld2006}. The literature on \ac{OS} also expands upon the formation of identities, particularly group identities by highlighting how humans understand \textit{self} and \textit{self-differentiation} \citep[pg.~749-750]{kinnvall2004} \citet[pg.~749-752]{kinnvall2004} poses that self-identification inherently requires a \textit{subject other}, and he debates the difference between categorization, both internally and externally, and membership regarding the role of societal groups in identity formation. Furthermore, \citet[pg.~751-752]{kinnvall2004} combines securitization with the subjectivity of \textit{other} to describe how both an actor's self-identity and subjective understanding of the \textit{other} are securitized to maintain a cohesive interpretation of society and the environment.

The way that identifies are essential to their securitization and to the securitization of the self. In \citepos{kinnvall2004} construction of \ac{OS}, the concept of \textit{abjection}---``unconscious self is neither an object nor a subject'' \citep[pg.~753]{kinnvall2004}---is a way to reduce uncertainty and increase \ac{OS}, both for individuals and collectives \citep[pg.~753-754]{kinnvall2004}. This unconscious state of abjection is a primary method for mitigating change, but it leads to the creation of collective identity ``when the familiar `stranger' is suddenly recognized as a threat'' \citep[pg.~753]{kinnvall2004}. This can be seen clearly in \ac{BiH} following the collapse of Yugoslavia and reemergence of ethnic identities that were weaponized to insight the ethnic conflicts in the 1990s. The dichotomy between self and other that leads to the creation of or membership to identities is particularly evident and destructive when the differences between self and other become ambiguous \citep[as cited in \citealt{kinnvall2004}, pg.~753]{norton1988}. 

This rather psychological approach to the formation of identity highlights the importance of maintaining \ac{OS} and reducing the ambiguities of the self. When actors are unsure of who they are or what is to come, they often turn to membership or creation of identities driven by fear or hatred of the \textit{other} rather than the simple dichotomous classification between self and ``stranger other'' \citep[pg.754~755]{kinnvall2004}. This formation of enemies is what leads to racism, anti-Semitism, ethnic conflicts, genocides, and conflict in general. This is not to say that ontological insecurity guarantees the formation of `enemies,' but it does increase the chances of these classifications forming.


\subsection{Routines \& Trust}
\label{subsec:rountine}

\ac{OS}'s goal of mitigating uncertainty and maintaining a consistent understanding of self and being relies largely on routinization and trust \citep[pg.~346]{mitzenOntologicalSecurityWorld2006}. An actor is never able to account for all possible outcomes in a situation, which is why \citet[pg.~36-37]{giddens1991} poses the concept of ``bracketing'' as a method that actors use to mitigate uncertainty, and this is formed by a trust in certain ``acceptable'' responses or reactions that are understood within a society. Bracketing includes virtually any form of knowledge, whether it is understood by society or just the actor, such as scientific facts, societal expectation, or identity, and a disruption to what is understood leads to ontological insecurity (\citealp[pg.~37-38]{giddens1991}; \citealp[pg.~346-347]{mitzenOntologicalSecurityWorld2006}). \citet[pg.~346]{mitzenOntologicalSecurityWorld2006} coins the collection of these ``bracketings'' as ``an actor's \textit{basic trust system}'', and this trust system is largely influenced by the various identities that an actor holds, providing the framework for how and why an actor has certain understood truths \citep[pg.~37-40]{giddens1991}.

This \textit{basic trust system} is the result of routinization and the unconscious decisions that result from routines \citep[pg.~346-347]{mitzenOntologicalSecurityWorld2006}. The mitigation of cognitive choice resultant of routines forms the braketings and trust systems used by actors to limit uncertainty, justify identities, and propagate \ac{OS} \citep[]{mitzenOntologicalSecurityWorld2006}. Routines also serve as a form of risk mitigation, both physical and ontological, as actors must always balance the possibility of risk through \textit{risk profiling} \citep[pg.~127-130]{giddens1991}. Despite this, attachment to routines may actively disregard risk out of fear or risk of altering or losing an identity \citep[pg.~347]{mitzenOntologicalSecurityWorld2006}. As \citet[pg.~350-351]{mitzenOntologicalSecurityWorld2006} claims, a ``rigid basic trust system'', heavily reliant on routines to mitigate disruption in self-understanding and identity, opens an actor to ontological insecurity due to the lack of cognitive distance from the routines themselves, which leads actors with ``rigid basic trust systems'' to act irrationally in order to preserve routine. This irrationality caused by routines and trust systems can be used to explain a variety of decisions that may detrimentally impact actors in order to maintain routines and the status quo.


\subsection{Narratives \& Collective Memory}
\label{subsec:narrative}

Narratives and collective memory are essential to the creation of national and ethnic identities, and they provide a framework for the collective understanding required of collective and institutionalized identities. States often control narratives through direct and indirect methods of discourse and action, political speech acts, memorialization of events, public holidays, legislature, and foreign policy \citep[pg.~615-617]{subotic2015}. National and ethnic narratives are constructed and maintained by ``narrative entrepreneurs''---political and intellectual elites, media, and education systems---and solidified through daily engagement with these narratives by the public \citep[pg.~615]{subotic2015}. Narratives act as a form of routine, as they provide meaning to state actions and form ``autobiographies'' for states that shape trust and collective memory \citep[pg.~614]{subotic2015}. When narratives are seen as ``rhetorical commonplace,'' they are difficult to dispute and become essential to the understanding of self, both of the state and its citizens \citep[pg.~615]{subotic2015}. This solidification of narratives ``brackets'' political debates and requires politicians to frame policy and discourse in alignment with narratives \citep[pg.~615-616]{subotic2015}.

%Add collective memory


\section{Ontological Security in BiH}
\label{sec:ont_BiH}

%find citations for referendum & history & parties
\ac{BiH} has a unique position regarding \ac{OS} as a state primarily constituent of three ethnicities, Muslim Bo\v{s}njaks, Catholic Croats, and Orthodox Serbs \citep[pg.~85]{Bringa_2000}. Following the collapse of the Socialist Federal Republic of Yugoslavia in, \ac{BiH} declared its independence in 1992, and was recognized by the European Community \citep[pg.~216]{davidPolicingMemoryBosnia2019}. Following independence, \ac{BiH} descended into conflict between the three ethnicities, Orthodox Serbs, Catholic Croats, and Muslim Bosniaks from April 1992 to October 1995, leading to immense damages, casualties, and displacement across \ac{BiH} \citep[pg.~216]{davidPolicingMemoryBosnia2019}. The conflict officially ended in December 1995 with the Dayton peace agreement, which recognized \ac{BiH} as an independent nation-state and defined the constitution that is still used in \ac{BiH} \citep[pg.~216]{davidPolicingMemoryBosnia2019}. The constitution provides a unique institutional framework that divides almost all functions of the government between the three major ethnicities \citep{oscemissiontobosniaandherzegovina1995}. As well dividing the institutions among the three ethnicities, two separate entities govern territory in \ac{BiH}---\ac{RS} which governs much of the Serbian ethnic areas in the north and west of \ac{BiH}, and \ac{FBiH} which governs much of the Croat and Bosniak ethnic areas in the west, south and center of \ac{BiH} \citep{oscemissiontobosniaandherzegovina1995}. Beyond the institutions of \ac{BiH}, three of the four major parties are also ethnically aligned with \ac{SNSD} being Serb aligned, \ac{HDZ} being Croat aligned, and \ac{SDA} being Bosniak aligned \citep{Parlamentarna_2026}. Not only are these parties ethnically aligned, but \ac{HDZ} is a sister party to HDZ in Croatia, and \ac{SDP}, which is the successor to the League of Communists of Yugoslavia, has sister parties in Croatia, Serbia, and North Macedonia. The complicated history and institutionalization of ethnic tensions leads to a unique case for \ac{OS} theory to address.

As highlighted above, \ac{OS} theory stresses that states seek to maintain stability in identity, routine, and trust through various means, but this becomes complicated in states like \ac{BiH} that have institutionalized divisions between ethnicities (See Section~\hyperref[sec:ontological_security]{\ref*{sec:ontological_security}}). The state seeks to maintain the \textit{Bosnian} identity, which encompasses all citizens and is what is projected externally to other states, but the various ethnicities also seek to maintain their respective ethnic identities, which may not always align with the state's projected identity. This complexity results in the opportunity to analyze how \ac{BiH} operationalizes \ac{OS} and the extent to which it is ontologically secure or insecure as a state. 

Identity is largely sustained through narratives, which often become engrained in collective memory as ``rhetorical commonplace,'' but actors are able to leverage these narratives through the \textit{activation} and \textit{deactivation} of narratives or parts of narratives, particularly in regard to political speech and policy \citep[pg.~615-616]{subotic2015}. This can be seen often in politics, and it is crucial for political discourse, debate, and policy. This can particularly be seen regarding events in the collective memory of \ac{BiH}, such as the war in the 1990s or World War II, leading to the possibility for collective memorialization of tragedy across all Bosnians or the weaponization of these memories for political gains \citep{davidPolicingMemoryBosnia2019}. The variation in narratives or parts of narratives have a major impact on the meaning of certain events, speech acts, and policy, and depending on how they are interpreted can support or hurt the understanding of state identity.


Part of Bosnian identity relies on 

Political speech acts are a primary way of understanding the operationalization of \ac{OS} theory at a state level because politicians have been given agency to speak for the social group that constitutes a state, and their speeches and actions are publicized to both internal and external actors \citep[pg.~280-281]{weldes1996}. It is ultimately politicians who speak and act on behalf of the state, and the interactions between politicians largely take place through parliamentary debate. It is the aim of this thesis to analyze how Bosnian politicians' speech acts in parliament shape the state's identity, routines, and trust, and if this aligns with \ac{OS} theory despite the complexity of Bosnian society and institutions.




%Unlike the other former Yugoslav states, there is a national narrative and identity that is separate from the ethnic narratives and identities, which inherently weakens the state's \ac{OS} \citep[pg.~216-217]{davidPolicingMemoryBosnia2019}

%The \textit{activation} and \textit{deactivation} of narrative aspects is the primary way that political actors frame political discourse and action, and the change if framing by political actors has the ability to change narratives overtime \citep[pg.~616]{subotic2015}. \ac{BiH} and other states with multiple ethnic identities have an additional level of nuance beyond the framing of the national narrative, the \textit{activation} and \textit{deactivation} of ethnic narratives. This additional layer allows for political discourse to \textit{activate} an ethnic narrative and \textit{deactivate} the national narrative or vice-versa. 

%Although states do not have a consciousness, political elites who act in the name of the state do, and the routinization of actions and meanings through narratives allows these actors to integrate these meanings into their own unconscious understanding of self (\citealp[pg.~614-615]{subotic2015}; \citealp[pg.~346-347]{mitzenOntologicalSecurityWorld2006}). Because of this internalization, the role of politicians in narrative creation is reflexive---politicians encourage consistent narratives that match their personal understanding of the narrative.


%The role of the state in these narratives is seen in all national discourses, but particularly in multi-ethnic, authoritarian states, who utilize the classification of ethnicities and a collective national identity to legitimize the power of the state (\citealp[pg.~85]{Bringa_2000}; \citealp[pg.~615-617]{subotic2015}). Despite the collective nature of national and ethnic identities, authoritarian regimes are able to suppress the ``narrative entrepreneurship'' of average citizens by controlling media outlets, education systems, and suppression of speech (\citealp[pg.~85]{Bringa_2000}; \citealp[pg.~615-617]{subotic2015}).




% Discuss the role of "Chosen Trauma" (Kinvall, 2004, pg.755)
% Combine with constructivism (Weldes)
% Begin wrapping up how statesmen shape national identity and \ac{OS} through speech acts


